\documentclass[11pt,a4paper]{article}

\usepackage[provide=*,spanish]{babel}
\usepackage[utf8]{inputenc}
\usepackage[T1]{fontenc}
\usepackage{lmodern}
\usepackage[margin=2.35cm]{geometry}
\usepackage{amsmath}
\usepackage{booktabs}
\usepackage{tabularx}
\usepackage{array}
\usepackage{enumitem}
\usepackage{listings}
\usepackage{xcolor}
\usepackage{microtype}
\usepackage{parskip}
\usepackage{titlesec}
\usepackage{hyperref}

\definecolor{azulWCD}{RGB}{17,67,112}
\definecolor{grisCodigo}{RGB}{247,247,247}
\definecolor{verdeCodigo}{RGB}{25,110,55}
\definecolor{rojoCodigo}{RGB}{145,35,35}
\definecolor{tituloColor}{RGB}{0,55,110}
\definecolor{reglaColor}{RGB}{0,90,160}

\hypersetup{
  colorlinks=true,
  linkcolor=azulWCD,
  urlcolor=azulWCD,
  citecolor=azulWCD,
  pdftitle={Implementación del RINDEX de la superficie Tyvek},
  pdfauthor={Jaime A. Betancourt}
}

\lstdefinestyle{cppstyle}{
  language=C++,
  backgroundcolor=\color{grisCodigo},
  basicstyle=\ttfamily\footnotesize,
  keywordstyle=\color{azulWCD}\bfseries,
  commentstyle=\color{verdeCodigo}\itshape,
  stringstyle=\color{rojoCodigo},
  numbers=left,
  numberstyle=\tiny\color{gray},
  numbersep=8pt,
  breaklines=true,
  frame=single,
  rulecolor=\color{gray!40},
  tabsize=2,
  showstringspaces=false,
  xleftmargin=1em,
  framexleftmargin=1em,
  columns=fullflexible,
  keepspaces=true
}
\lstset{style=cppstyle}

\titleformat{\section}{\normalfont\Large\bfseries\color{azulWCD}}
  {\thesection.}{0.6em}{}
\titleformat{\subsection}{\normalfont\large\bfseries}
  {\thesubsection.}{0.5em}{}
\setlist{itemsep=0.25em,topsep=0.4em}
\setlength{\emergencystretch}{3em}
\renewcommand{\arraystretch}{1.18}

\newcommand{\codigo}[1]{\nolinkurl{#1}}
\newcommand{\archivo}[1]{\path{#1}}

\begin{document}
\pagestyle{empty}

%------------------------------------------------
% PORTADA
%------------------------------------------------

\begin{titlepage}
\begin{center}

\vspace*{0.5cm}
\textcolor{reglaColor}{\rule{\linewidth}{1.2pt}}

\vspace{1.3cm}
{\Large\scshape Universidad Industrial de Santander}\\[0.2cm]
{\large Escuela de Física --- Facultad de Ciencias}\\[0.2cm]
{\large Doctorado en Física}

\vspace{2.2cm}
{\Huge\bfseries\color{tituloColor} Implementación del índice\\[0.2cm]
de refracción en la superficie\\[0.2cm] óptica del Tyvek}

\vspace{1.8cm}
{\large\textbf{Informe técnico}}

\vspace{0.6cm}
\begin{minipage}{0.82\linewidth}
\centering
\textit{``Detector Cherenkov de agua con medios agua--NaCl''}
\end{minipage}

\vfill

{\large\textbf{Autor}}\\[0.15cm]
Jaime A. Betancourt\\[0.1cm]
Doctorado en Física, Universidad Industrial de Santander

\vspace{1.2cm}
\textcolor{reglaColor}{\rule{0.5\linewidth}{0.6pt}}

\vspace{1cm}
11 de septiembre de 2026

\vspace{1cm}
\textcolor{reglaColor}{\rule{\linewidth}{1.2pt}}

\end{center}
\end{titlepage}

\pagestyle{plain}

\begin{abstract}
Este informe documenta la corrección introducida en la representación óptica
del recubrimiento interno de Tyvek del detector Cherenkov de agua. La versión
revisada incorpora la propiedad \codigo{RINDEX} en la tabla
\codigo{linerOpticalPT}, asociada a una superficie
\codigo{dielectric_dielectric} con acabado \codigo{groundbackpainted}. También
se incorporan verificaciones en \codigo{SaltyWCD} para comprobar que tanto el
medio activo como la superficie reflectora contienen las propiedades mínimas
requeridas antes de construir la geometría. La corrección mantiene una única
descripción intrínseca del Tyvek para agua pura y las soluciones con 2.5, 5.0 y
10.0\,\% de NaCl; la dependencia con la concentración entra mediante las
propiedades ópticas del medio activo.
\end{abstract}

\tableofcontents

\section{Objeto de la modificación}

El WCD utiliza una lámina de Tyvek como recubrimiento interno difusamente
reflector. En la implementación analizada, el Tyvek estaba representado en dos
niveles:

\begin{enumerate}
  \item como un material volumétrico de HDPE, denominado
  \codigo{Tyvek_HDPE};
  \item como una superficie óptica, denominada
  \codigo{TyvekOpticalSurface}, aplicada a las fronteras entre el medio activo
  y las láminas superior, inferior y lateral.
\end{enumerate}

El material volumétrico ya disponía de índice de refracción y longitud de
absorción. Sin embargo, la tabla propia de la superficie
\codigo{linerOpticalPT} no contenía \codigo{RINDEX}. Esta omisión es relevante
porque la superficie se configuró con acabado \codigo{groundbackpainted}; en
este modelo, el proceso de frontera óptica necesita el índice de refracción de
la capa superficial.

\section{Representación volumétrica del Tyvek}

El material se conserva como una aproximación de polietileno de alta densidad:

\begin{lstlisting}
HDPE = new G4Material("Tyvek_HDPE", 0.94 * g/cm3, 2);
HDPE->AddElement(elC, 2);
HDPE->AddElement(elH, 4);
\end{lstlisting}

Su tabla volumétrica \codigo{linerPT1} contiene:

\begin{lstlisting}
linerPT1->AddProperty("RINDEX",
                      water1PhotonEnergy,
                      linerRefIndex,
                      water1ArrEntries);

linerPT1->AddProperty("ABSLENGTH",
                      water1PhotonEnergy,
                      linerAbsLen,
                      water1ArrEntries);

HDPE->SetMaterialPropertiesTable(linerPT1);
\end{lstlisting}

El índice utilizado es $n=1.49$ y la longitud de absorción es
$0.1\,\mathrm{mm}$ en los puntos de la malla. Estas propiedades describen el
transporte de un fotón que ingresa físicamente al volumen de HDPE.

\section{Representación superficial}

La interacción de un fotón en la frontera medio activo--Tyvek se representa
mediante \codigo{LinerOptSurf}. Antes de la corrección,
\codigo{linerOpticalPT} contenía reflectividad y parámetros del modelo
\textit{UNIFIED}, pero carecía del índice de refracción de la capa pintada.

\subsection{Corrección implementada}

Se añadió:

\begin{lstlisting}
linerOpticalPT->AddProperty("RINDEX",
                            water1PhotonEnergy,
                            linerRefIndex,
                            water1ArrEntries);
\end{lstlisting}

El bloque superficial queda conceptualmente compuesto por:

\begin{lstlisting}
linerOpticalPT = new G4MaterialPropertiesTable();

linerOpticalPT->AddProperty(
    "RINDEX",
    water1PhotonEnergy,
    linerRefIndex,
    water1ArrEntries);

linerOpticalPT->AddProperty(
    "REFLECTIVITY",
    water1PhotonEnergy,
    linerReflectivity,
    water1ArrEntries);

linerOpticalPT->AddProperty(
    "BACKSCATTERCONSTANT",
    water1PhotonEnergy,
    linerBackScatter,
    water1ArrEntries);
\end{lstlisting}

También se conservan \codigo{SPECULARLOBECONSTANT} y
\codigo{SPECULARSPIKECONSTANT}. Finalmente, la tabla se asigna mediante:

\begin{lstlisting}
LinerOptSurf->SetModel(unified);
LinerOptSurf->SetType(dielectric_dielectric);
LinerOptSurf->SetFinish(groundbackpainted);
LinerOptSurf->SetSigmaAlpha(fSigmaAlpha);
LinerOptSurf->SetMaterialPropertiesTable(linerOpticalPT);
\end{lstlisting}

\subsection{Parámetros del modelo UNIFIED}

\begin{center}
\begin{tabularx}{\textwidth}{>{\ttfamily}l X}
\toprule
\normalfont\textbf{Propiedad} & \textbf{Interpretación}\\
\midrule
RINDEX & Índice de refracción de la capa óptica superficial.\\
REFLECTIVITY & Probabilidad espectral asociada a la reflexión del recubrimiento.\\
BACKSCATTERCONSTANT & Fracción correspondiente a retrodispersión.\\
SPECULARLOBECONSTANT & Reflexión especular respecto a la normal de una microfaceta.\\
SPECULARSPIKECONSTANT & Reflexión respecto a la normal macroscópica.\\
sigmaAlpha & Anchura angular de la distribución de normales de microfacetas.\\
\bottomrule
\end{tabularx}
\end{center}

Con los valores utilizados, las componentes explícitas son 0.02 para el lóbulo
especular, 0 para el pico especular y 0.02 para retrodispersión. La componente
difusa restante es calculada por el modelo:

\begin{equation}
P_{\text{difusa}}
=1-P_{\text{lóbulo}}-P_{\text{pico}}-P_{\text{retro}}
=0.96.
\end{equation}

\section{Dependencia con el medio detector}

No se crean cuatro superficies Tyvek diferentes. El recubrimiento mantiene una
única curva intrínseca de índice, reflectividad y rugosidad. La dependencia con
la concentración de NaCl se introduce mediante el índice espectral del medio
activo seleccionado:

\begin{center}
\begin{tabular}{lc}
\toprule
Medio activo & $n(589.3\,\mathrm{nm})$\\
\midrule
Agua pura & 1.3330\\
Agua + 2.5\,\% NaCl & 1.3397\\
Agua + 5.0\,\% NaCl & 1.3436\\
Agua + 10.0\,\% NaCl & 1.3594\\
\bottomrule
\end{tabular}
\end{center}

Geant4 combina el índice del medio incidente con la descripción de la
superficie. Por ello pueden cambiar las probabilidades asociadas a reflexión de
Fresnel, refracción y reflexión interna total, aunque la curva propia del Tyvek
sea la misma.

Esta aproximación es apropiada mientras no se disponga de mediciones de
reflectividad del sistema Tyvek--solución salina para cada concentración.

\section{Invocación desde SaltyWCD}

\codigo{SaltyWCD} construye tres volúmenes de Tyvek:

\begin{itemize}
  \item \codigo{logTyvekTop};
  \item \codigo{logTyvekBot};
  \item \codigo{logTyvekSide}.
\end{itemize}

Todos utilizan \codigo{mat.HDPE}. Después de posicionarlos, se crean tres
superficies de frontera ordenadas desde el volumen activo hacia el Tyvek:

\begin{lstlisting}
new G4LogicalBorderSurface("WaterTyvekTopSurface",
                           physTank,
                           physTyvekTop,
                           mat.LinerOptSurf);

new G4LogicalBorderSurface("WaterTyvekBottomSurface",
                           physTank,
                           physTyvekBot,
                           mat.LinerOptSurf);

new G4LogicalBorderSurface("WaterTyvekSideSurface",
                           physTank,
                           physTyvekSide,
                           mat.LinerOptSurf);
\end{lstlisting}

La geometría no necesita redefinir \codigo{RINDEX}: recibe
\codigo{LinerOptSurf} ya configurada en \codigo{Materials}.

\section{Verificaciones añadidas}

Se reforzó \codigo{SaltyWCD::BuildDetector} para verificar la tabla del medio:

\begin{lstlisting}
if (!opticalTable ||
    !opticalTable->GetProperty("RINDEX") ||
    !opticalTable->GetProperty("ABSLENGTH")) {
    G4Exception("SaltyWCD::BuildDetector",
                "WCD-OPT-001",
                FatalException,
                "The selected WCD medium requires "
                "RINDEX and ABSLENGTH.");
}
\end{lstlisting}

También se comprueba la tabla superficial:

\begin{lstlisting}
if (!tyvekSurfaceTable ||
    !tyvekSurfaceTable->GetProperty("RINDEX") ||
    !tyvekSurfaceTable->GetProperty("REFLECTIVITY")) {
    G4Exception("SaltyWCD::BuildDetector",
                "WCD-OPT-002",
                FatalException,
                "The Tyvek groundbackpainted surface "
                "requires RINDEX and REFLECTIVITY.");
}
\end{lstlisting}

Estas excepciones impiden iniciar una campaña con una configuración óptica
incompleta.

\section{Alcance y limitaciones}

La corrección:

\begin{itemize}
  \item no modifica la composición ni la densidad del HDPE;
  \item no altera las matrices ópticas del agua o de las soluciones salinas;
  \item no modifica la reflectividad espectral disponible;
  \item no crea una dependencia empírica de la reflectividad con la salinidad;
  \item no cambia el Pyrex, el PMT ni el acero;
  \item no modifica los procesos ópticos de la lista de física.
\end{itemize}

Las superficies \codigo{G4LogicalBorderSurface} son direccionales. La
implementación cubre explícitamente los fotones que viajan desde
\codigo{physTank} hacia las láminas de Tyvek. Si una validación demuestra que
los fotones entran al volumen de HDPE y regresan, deberá evaluarse la creación
de las superficies inversas.

\section{Validación recomendada}

La revisión estática confirmó que:

\begin{itemize}
  \item las cinco matrices espectrales consideradas contienen 30 valores:
  \codigo{water1PhotonEnergy}, \codigo{linerRefIndex},
  \codigo{linerReflectivity}, \codigo{linerBackScatter} y
  \codigo{linerAbsLen};
  \item \codigo{RINDEX} y \codigo{REFLECTIVITY} comparten la misma malla;
  \item las tablas volumétrica y superficial permanecen separadas;
  \item las tres fronteras del tanque emplean \codigo{LinerOptSurf}.
\end{itemize}

Permanecen pendientes:

\begin{enumerate}
  \item compilar el proyecto con MEIGA y Geant4 10.07.p04;
  \item ejecutar fotones ópticos monoenergéticos hacia las paredes;
  \item contar reflexión, absorción y transmisión;
  \item verificar el cierre
  \begin{equation}
  \eta_{\mathrm{Refle}}+
  \eta_{\mathrm{Abs}}+
  \eta_{\mathrm{Trans}}+
  \eta_{\mathrm{Otros}}\simeq 1;
  \end{equation}
  \item repetir la prueba para agua pura y 2.5, 5.0 y 10.0\,\% de NaCl.
\end{enumerate}

\section{Archivos corregidos}

\begin{itemize}
  \item \archivo{Materials-mejorado-optica-Tyvek-RINDEX.cc};
  \item \archivo{SaltyWCD-mejorado-optica-Tyvek-RINDEX.cc}.
\end{itemize}

Para integrarlos en el repositorio deberán compararse con la rama activa y
renombrarse como \archivo{Materials.cc} y \archivo{SaltyWCD.cc}. La revisión
realizada no modifica ni publica el repositorio externo.

\section{Conclusiones}

La incorporación de \codigo{RINDEX} en \codigo{linerOpticalPT} completa la
información mínima esperada para la superficie
\codigo{dielectric_dielectric/groundbackpainted}. La distinción entre la tabla
volumétrica del HDPE y la tabla de la superficie evita mezclar transporte en el
material con interacción en la frontera. Las verificaciones añadidas en
\codigo{SaltyWCD} convierten una posible pérdida silenciosa de fotones en un
error explícito y trazable.

\end{document}
